% 第 35 章　把光当成波
\chapter{把光当成波}

\step{反常识开场}

\begin{mainline}
吹一个肥皂泡，盯着它看。肥皂膜明明是透明的，阳光明明是无色的，可泡的表面却滚动着一圈圈鲜艳的颜色，红、绿、蓝、紫像液体一样流动。这些颜色没有画上去，没有染上去，一阵风吹过，花纹立刻换一副面孔。颜色到底藏在哪里？

再看一张废旧光盘。把它倾斜着对准灯光，盘面会炸开一道完整的彩虹——可是光盘表面是光亮的塑料，没有一滴颜料。而一块同样光亮的玻璃板，或者一面镜子，就绝对变不出这道彩虹。光盘比镜子多出来的东西是什么？

这两件事还只是热身。真正让人坐不住的是第三个现象。让一束很细的光穿过挡板上两条紧挨着的细缝，照到后面的墙上。按常识，墙上应该出现两条亮线，中间暗，两边暗，仅此而已。实际看到的却是一整排等间距的亮暗条纹，像斑马一样铺开十几条。更离奇的还在后面：找准一条\emph{暗}条纹的正中央——那里本来一点光都没有——然后伸手挡住两条缝中的一条。光变少了，墙上总亮度变小了，可那个原本全黑的位置，竟然\emph{亮了起来}。

挡住一部分光，黑暗的地方反而亮了。这一章所有的概念，都是被这一句反常识的话逼出来的。
\end{mainline}

\step{先猜一猜}

老规矩，先把你的直觉押在桌上，不要偷看后面的答案。

第一题。夜里看一盏远处的路灯（或者把手机手电筒开到最亮，放在房间另一头），把两根手指并在一起，中间留一条很细的缝，举到眼前透过缝去看。你会看到什么？（a）一条被挤窄的亮光；（b）几道沿着手指方向的亮纹；（c）一列\emph{垂直于}手指方向排开的明暗条纹。注意选项（c）的方向——如果你直觉选的不是它，请记住这份别扭。

第二题。双缝实验里，墙上有一条完全黑暗的条纹。现在挡住两条缝中的一条，让进光量减少一半。那个黑暗位置会：（a）变得更黑；（b）不变，还是黑的；（c）变亮。

第三题。光盘上的彩虹，是因为光盘塑料里含有某种会随角度变色的物质吗？如果把光盘表面那层彩虹色的结构磨掉，磨出一个彻底光滑的圆盘，彩虹还在不在？

写下你的三个答案和理由。本章结尾我们逐一对账。

\step{动手实验}

\begin{experiment}[指缝里的条纹与光盘上的彩虹]
\textbf{材料}：一部带手电筒功能的手机、一张废旧 CD 或 DVD、一间能变暗的房间。选做：洗洁精、水、一根弯成环的铁丝。

\textbf{步骤一}：关掉房间的灯，把手机手电筒放在两三米外，对着你。伸出一只手，食指和中指并拢，再稍稍用力，让两指之间只留下一条头发丝般的细缝。把缝举到眼前，透过它看那盏“灯”。眯起眼睛仔细看：你看到的不是一条细光，而是沿着\emph{垂直于指缝}的方向排开的一小列明暗相间的条纹，缝捏得越紧（越窄），条纹分得越开。光线的直觉说“缝窄了，看到的东西应该变窄”，实验偏要反着来。

\textbf{步骤二}：拿起光盘，让手电筒的光斜照在有字的反面（光亮面），转动光盘，观察反射光。在某个角度你会看到盘面铺展开一道彩虹；换一个角度，彩虹换一种排列。再用同样的动作照一照手边的镜子或玻璃——什么颜色也没有。

\textbf{步骤三（选做）}：洗洁精兑水，用铁丝环蘸起一层皂膜，竖直拿着，对着灯光看膜的反射光。几分钟后膜上的颜色会分层流动，出现水平的彩色带；如果耐心看到最后，膜的顶部会出现一块\emph{发黑}的区域——那不是膜破了，它往往出现在破裂之前。

\textbf{记录}：指缝条纹沿哪个方向排开？缝越窄条纹越挤还是越开？光盘的彩虹和镜子的反光有什么区别？皂膜顶部为什么先变黑？把观察写下来，本章会逐个拆开。
\end{experiment}

还有一个零成本的观察：雨后的马路上，一滩积水表面漂着薄薄一层油花时，油花会显出斑斓的彩色。同样是透明的水、透明的油、无色的阳光，颜色却出现了。它和肥皂泡是同一个秘密。

\step{旧模型遇到麻烦}

第 33 章和第 34 章，我们把光当成射线：它走直线，碰到镜面按反射定律弹开，穿过界面按折射定律拐弯，透镜把它聚成清晰的像。这套“光线模型”战果辉煌——相机、望远镜、眼镜、光纤，全都靠它设计出来。它那么好使，为什么要自找麻烦？

因为麻烦自己找上门了，而且一来就是一串。

麻烦一：影子不肯利落。光线模型说光走直线，影子应该有刀切般的边缘：亮就是亮，暗就是暗。可你仔细看阳光下任何影子的边缘，都有一条模糊的过渡带；用小孔成像时（第 34 章结尾的伏笔），孔缩得太小，像不但没变得更锐利，反而整个糊掉了。缩小孔径在光线模型里只会“减少进光量”，没有理由改变图案的形状——可形状确实变了。

麻烦二：方向对不上。指缝实验里，指缝是竖直的，光线模型预言你看到的光应该被“挤”在竖直方向上；实际看到的条纹却沿着\emph{水平}方向排开，缝在哪个方向捏窄，图案就往垂直的方向摊开。直线传播完全无法解释光为什么往“不该去的方向”跑。

麻烦三：颜色没地方放。光线模型里，镜面反射只把光送到一个确定的方向；光盘表面无非是许多镜面，凭什么把白光拆成彩虹，还随角度变化？颜料说不通——彩虹转动光盘就换位置，物质不会跟着角度换成分。

麻烦四，也是最致命的一条：双缝。两束光在墙上同一地点相遇，光线模型只会做一种算术——两束光的亮度相加，只会更亮或不变，绝不可能更暗。可实验里偏偏出现了稳定的、完全黑暗的条纹；而挡住其中一束光之后，黑暗处竟然亮了。“光加光等于暗”，在光线模型的语言里根本写不出这句话，就像算式 \(1+1=0\) 一样刺眼。

四桩麻烦有一个共同点：出事的场合里，障碍物的尺寸（指缝、小孔、光盘的刻痕）都很\emph{小}。光线模型在大尺度上无往不利，一到小尺度就节节败退。这不是“光线模型算错了一点”，而是它漏掉了光的一种本性。要补上这个本性，我们得回到第 17 章学过的老朋友——波。

\step{发明新概念}

\begin{mainline}
\textbf{第一个新概念：光是一种波，有波长。}第 25 章已经告诉我们，光是电磁波，在真空里以速度 \(c\) 前进。波有两个基本身份：频率（每秒振动多少次）和波长（一个完整的起伏在空间占多长）。可见光的波长大约在 \(400\) 纳米（紫光）到 \(700\) 纳米（红光）之间——一根头发丝的直径，大约等于一百多个可见光波长排起来。光的全部颜色，就是波长的花名册：波长偏长显红，偏短显紫。平时我们察觉不到光在“波动”，原因很简单：它的波长太短了，短到日常尺度上它表现得像一条直线。只有障碍物小到和波长差不多（或者大不了几百倍）时，波的本性才藏不住。

\textbf{第二个新概念：惠更斯原理——波前上每一点都是新的波源。}波传到某处，那个波前（同一时刻振动走到同一节拍的点连成的面）上的每一点，都可以当作一个小小的子波源，各自发出向前的球面子波；下一时刻的波前，就是所有这些子波叠加出来的包络。它像什么？像体育场里的人浪：没有人从头到尾跑完一圈，每个观众只是被旁边的人带着起立坐下，可“人浪”这个波形自己走完了全场。它不像什么？不像在波前上真的撒了一排小灯泡——子波不是独立的实体，而是我们计算“波接下来会走到哪”的一种方法。

先别急着用它解释怪事，它首先得解释“平常”：波为什么通常走直线？想象一列宽阔的平面波前，上面每一点都发出球面子波。侧向的子波来自密密麻麻的相邻各点，彼此节拍参差，叠加起来相互抵消；只有在正前方，所有子波同步到达，叠出一道新的平直波前——包络面依然平直，波就直着走。日常尺度下，沿直线前进的正是这张“子波包络”，第 33 章的“光线”不过是垂直于包络面的那条参考线。反射定律和折射定律也能从子波的几何里推出来，这里不展开；重点是，惠更斯原理没有把旧模型推翻，而是把旧模型解释成了波在“宽阔波前”极限下的样子。

但这条原理真正的锋芒在小尺度。它一句话就解释了光为什么会绕弯：波前擦过障碍物边缘时，边缘附近那些没被挡住的点照常发出子波，子波是球面的，自然会有一部分拐进几何阴影里去。障碍物越小，被“放行”的波前越窄，这点残存的波前就越像一个点波源，把球面波撒向所有方向——这就是衍射。光并没有在边缘上“反弹拐弯”，是直线传播的近似在小尺度下到期了。

\textbf{第三个新概念：相位与叠加——波按节拍做加减法。}描述一个波，除了“振多强”（振幅），还要说“此刻振到哪一步”（相位）。两个波在同一点相遇时，该点的振动是两个振动之和：波峰遇上波峰，同拍相加，振动加倍；波峰遇上波谷，反拍相消，互相抵消。这件事你在第 18 章已经听过了：两只播放同一首歌的音箱之间，存在几个声音特别弱的“安静点”——两只音箱发出的声波到那些点的路程不同，到达时恰好反拍。现在把同一件事搬到光上：光波也是振动（振动的是电场和磁场，而不是某种介质的位移——这是它和水波、声波最不一样的地方，不需要任何“水面”来承载），所以两列光波相遇，同样按节拍做加减。同拍相加的地方就是亮条纹，反拍相消的地方就是暗条纹。那句光线模型说不出的“光加光等于暗”，用波的语言只是半句话：波峰加波谷。
\end{mainline}

\begin{center}
\begin{tikzpicture}[scale=1.0]
  % 挡板与双缝
  \draw[very thick] (-2.6,-1.8) -- (-2.6,-0.35);
  \draw[very thick] (-2.6,-0.15) -- (-2.6,0.15);
  \draw[very thick] (-2.6,0.35) -- (-2.6,1.8);
  \node[left] at (-2.6,-0.25) {缝};
  % 入射波前
  \foreach \x in {-4.6,-4.0,-3.4} {\draw[blue!60!black] (\x,-1.5) -- (\x,1.5);}
  \draw[-{Stealth[length=2.5mm]},blue!60!black] (-4.6,0) -- (-3.0,0);
  \node[blue!60!black,above] at (-3.9,1.5) {入射平面波};
  % 两缝发出的子波（圆弧）
  \foreach \r in {0.5,1.0,1.5,2.0} {
    \draw[gray!60,dashed] (-2.6,0.25) ++(0:\r) arc (0:55:\r);
    \draw[gray!60,dashed] (-2.6,0.25) ++(0:\r) arc (0:-55:\r);
    \draw[gray!60,dashed] (-2.6,-0.25) ++(0:\r) arc (0:55:\r);
    \draw[gray!60,dashed] (-2.6,-0.25) ++(0:\r) arc (0:-55:\r);}
  % 屏幕与条纹
  \draw[very thick] (2.6,-1.8) -- (2.6,1.8);
  \foreach \y/\w in {-1.4/0.5,-0.9/0.9,-0.45/0.5,0/0.9,0.45/0.5,0.9/0.9,1.4/0.5}
    {\fill[red!60!black] (2.66,\y-0.09) rectangle (2.78,\y+0.09);}
  \node[below] at (2.7,-1.8) {屏幕上的亮条纹};
  % 到屏上一点的两条路径
  \draw[red!70!black,thick] (-2.6,0.25) -- (2.6,0.9);
  \draw[red!70!black,thick] (-2.6,-0.25) -- (2.6,0.9);
  \node[red!70!black,right] at (0.4,0.75) {两条路，长度不同};
\end{tikzpicture}
\end{center}

\begin{mainline}
\textbf{第四个新概念：光程差是明暗的裁判。}双缝实验里，两条缝被同一个波前照亮，可以作为两个同拍的子波源。墙上某一点到两条缝的距离一般不相等：波从两条缝出发时同拍，走了不同的路，到达时的节拍差就完全由“路程差”决定。路程差正好是波长的整数倍，两列波到达时峰对峰——亮；路程差是波长的一半、一倍半、两倍半……峰对谷——暗。墙上那排斑马纹，不过是一张“路程差分布图”。这也顺便解释了为什么两条缝必须靠得很近、必须被同一束光照亮：路程差要小到只有几个波长，节拍差才稳定、才看得出条纹。

\textbf{第五个新概念：相干性——不是任何两束光都能稳定干涉。}那为什么用两支手电筒照墙，永远照不出暗条纹？因为“反拍相消”要维持得够久才看得见。普通光源（太阳、灯泡、手电筒）发出的光，是一截一截毫无默契的短波列，两束光之间的节拍关系每秒钟要随机翻转亿亿次，亮暗条纹的位置跟着疯狂闪烁，眼睛和相机只能拍到糊成一团的平均——均匀亮度。两个波源的节拍关系稳定，称为\emph{相干}。双缝实验的聪明之处，是让\emph{同一}波前分成两路：不管波列怎么随机翻转，两路翻得一模一样，节拍差原地不动，条纹就站稳了。激光笔之所以能做家庭干涉实验，是因为激光的波列远比普通光源整齐，相干时间长得多。所以“两支手电筒照不出条纹”不是仪器不够好，是那两束光之间根本没有可供稳定的节拍关系。
\end{mainline}

\step{一句话模型}

\begin{oneline}
把光当成波：屏幕上每一点的明暗，由“沿不同路径到达该点的振动，此刻是同步还是反步”决定——路程差凑整波长就加亮，差半个波长就抵消变暗；障碍物的尺寸一接近波长，光就不再走直线，而是按惠更斯原理重新铺开。
\end{oneline}

\step{数学翻译器}

\begin{mainline}
先把双缝条纹翻译成式子。设两条缝相距 \(d\)，屏幕在远处 \(L\)，考察屏上与正中央成角度 \(\theta\) 的一点。从图上的几何关系可以读出，两条路径的路程差为
\[
  \delta = d\sin\theta .
\]
上一步说过裁判规则：路程差凑整波长则亮，差半个波长则暗。所以
\[
  d\sin\theta = m\lambda \ \text{（亮条纹）},\qquad
  d\sin\theta = \left(m+\tfrac12\right)\lambda \ \text{（暗条纹）},\qquad m=0,1,2,\dots
\]
读一遍每一项：左边是“两条路差了多少”，右边是“差了几个波长”；\(\theta\) 回答“条纹在屏幕的哪个方向”，\(\lambda\) 是光自己的尺子。角度很小时，相邻亮纹在屏上的间距为
\[
  \Delta y \approx \frac{L\lambda}{d}.
\]
立刻做三道例行检查。第一，\(\theta=0\)（正中央）：\(\delta=0\)，任何波长都满足亮纹条件——所以用白光做实验，中央条纹永远是白色，所有颜色在这里同拍。第二，令 \(d\) 变小：\(\Delta y\) 变大，条纹变稀；\(d\) 趋向零时条纹间距趋向无穷大，整个屏幕均匀变亮——两条缝合并成一个光源，干涉消失，合理。第三，令 \(\lambda\) 趋向零（相对 \(d\) 而言）：\(\Delta y\) 趋向零，条纹密到眼睛分不开，屏幕上只剩两条缝的几何影子——波动模型自动退化回第 33 章的光线模型。新模型把旧模型收编为极限情形，这是它合格的第一张证书。

来一组真实数字。红色激光笔的波长约 \(650\) 纳米，用针在不透光纸上扎一对相距约 \(0.25\) 毫米的小孔，屏幕放在 \(1\) 米外：
\[
  \Delta y \approx \frac{1\ \text{米}\times 650\times10^{-9}\ \text{米}}{0.25\times10^{-3}\ \text{米}}
  \approx 2.6\ \text{毫米}.
\]
毫米级的条纹间距，肉眼刚好能数——这就是为什么这个实验在黑暗的客厅里就能做成。反过来看，如果用 \(d=5\) 毫米的“宽缝对”，条纹间距只有约 \(0.13\) 毫米，密得糊成一片。干涉条纹从来都在，只是大多数时候密得看不见。

第二条式子给单缝衍射。一条宽为 \(a\) 的缝，按惠更斯原理可看作一排子波源；把缝两端的子波比较，当它们的路程差达到一个波长时，缝内各对子波恰好两两抵消，出现第一条暗纹：
\[
  a\sin\theta = \lambda \quad\text{（第一暗纹）},
\]
中央亮纹的角宽度约为 \(2\lambda/a\)。例行检查：\(a\) 越大，\(\theta\) 越小——宽缝几乎不衍射，光直着走，回到光线模型；\(a=\lambda\) 时 \(\sin\theta=1\)，\(\theta=90^\circ\)，第一暗纹被推到天边，光均匀撒向整个前方——孔小到和波长一样，就成了一个球面波源；\(a\) 介于其间，缝越窄图案越宽。这就是指缝实验里“捏得越紧，条纹越开”的全部原因，也解释了条纹为什么沿\emph{垂直于缝}的方向铺开：衍射只发生在“缝被捏窄的那个方向”上。
\end{mainline}

\begin{center}
\begin{tikzpicture}[scale=1.0]
  % 单缝衍射强度分布
  \draw[-{Stealth[length=2.5mm]}] (-4.2,0) -- (4.4,0) node[below left] {屏上位置};
  \draw[-{Stealth[length=2.5mm]}] (0,0) -- (0,2.4) node[below left] {亮度};
  \draw[very thick,blue!70!black,smooth] plot coordinates {
    (-4.0,0.00) (-3.6,0.03) (-3.2,0.00) (-2.9,0.06) (-2.6,0.00)
    (-2.2,0.15) (-1.9,0.00) (-1.4,0.50) (-0.8,1.40) (0,2.00)
    (0.8,1.40) (1.4,0.50) (1.9,0.00) (2.2,0.15) (2.6,0.00)
    (2.9,0.06) (3.2,0.00) (3.6,0.03) (4.0,0.00)};
  \draw[dashed] (0,0) -- (0,2.0);
  \draw[dashed] (1.9,0) -- (1.9,0.25);
  \draw[dashed] (-1.9,0) -- (-1.9,0.25);
  \node[below] at (1.9,0) {第一暗纹};
  \node[below] at (-1.9,0) {第一暗纹};
  \draw[{Stealth[length=2mm]}-{Stealth[length=2mm]},red!70!black] (-1.9,1.0) -- (1.9,1.0);
  \node[red!70!black,above] at (0,1.0) {中央亮纹宽度 \(\approx 2\lambda/a\)};
\end{tikzpicture}
\end{center}

\begin{mainline}
再估一次头发丝。把一根头发（直径约 \(70\) 微米）放在激光笔前，头发后面 \(1\) 米处的屏上，阴影里会嵌着一条亮线，两侧铺开衍射条纹，第一暗纹的方向角约
\[
  \theta \approx \frac{\lambda}{a} = \frac{550\times10^{-9}\ \text{米}}{70\times10^{-6}\ \text{米}}
  \approx 0.008\ \text{弧度},
\]
在 \(1\) 米外的屏上相当于离中心约 \(8\) 毫米。一根头发就能当一台波长测量仪：量出条纹位置，反推波长，数量级就是几百纳米。

第三条式子给光栅。光盘上那圈螺旋刻痕，相邻轨道相距约 \(1.6\) 微米——这是一把刻着几千条“缝”的尺子，叫衍射光栅。大量等间距的缝让干涉条纹被极度锐化，亮纹方向满足
\[
  d\sin\theta = m\lambda ,
\]
形式与双缝一模一样，区别只在亮纹变得又细又亮。代入真实数字：\(d=1.6\) 微米，第一级（\(m=1\)）亮纹处，紫光（\(400\) 纳米）\(\theta\approx 14^\circ\)，红光（\(650\) 纳米）\(\theta\approx 24^\circ\)——同一级条纹里，不同颜色被掰向不同角度，白光就这样被拆成彩虹。镜子没有刻痕，没有这把尺子，自然什么颜色也变不出来。你猜对第三题了吗？彩虹不在塑料里，在刻痕的间距里；磨平刻痕，彩虹就死了。

“路程差决定明暗”这把钥匙不只管光。降噪耳机是它在声学里的工业版：耳机上的麦克风先“听”到外界噪声，芯片立刻造出一个反拍的声波送进你耳朵，噪声与反噪声在鼓膜处相消——“声音减声音等于安静”，和“光减光等于暗”是同一条原理。区别只在光的波长太短，想给光造一副“降噪耳机”，就得让两路光的光程差控制在一个波长以内，这正是干涉仪器（比如测引力波的激光干涉仪）难造的原因。
\end{mainline}

\begin{mathbridge}
“同拍相加、反拍相消”可以用相位写得更干净。路程差 \(\delta\) 对应的相位差为
\[
  \Delta\varphi = \frac{2\pi}{\lambda}\,\delta ,
\]
即每差一个波长，节拍差一整圈 \(2\pi\)。两列等振幅的波叠加，合成振幅正比于 \(\cos(\Delta\varphi/2)\)，而亮度（光强）正比于振幅的平方：
\[
  I = I_0\cos^2\!\frac{\Delta\varphi}{2}.
\]
检查：\(\Delta\varphi=0\)，\(I=I_0\)，最亮；\(\Delta\varphi=\pi\)，\(\cos(\pi/2)=0\)，全暗；在这之间平滑过渡——所以真实条纹是渐变的，不是刀切的亮暗。能量守恒也藏在这行式子里：把 \(I\) 对整个屏幕平均，结果恰好等于两束光各自亮度的总和，暗处“消失”的能量，一分不少地堆进了亮处。

薄膜干涉的算法是同一把钥匙，只是“路程”要换成“光程”。光在折射率为 \(n\) 的膜里走厚度 \(t\)，由于膜中波长被压缩为 \(\lambda/n\)，节拍按 \(nt\) 计算，上下表面两束反射光的光程差为 \(2nt\)。另有一个细节：光从光疏介质射向光密介质反射时，相位会额外翻转半个波长（半波损失），肥皂膜只有上表面的反射带这个翻转。于是反射相消的条件是
\[
  2nt = m\lambda .
\]
相机镜头和眼镜片上的增透膜就按这条式子制造：取膜的厚度使 \(2nt=\lambda/2\) 附近相消，即 \(t\approx \lambda/(4n)\)。用折射率约 \(1.38\) 的氟化镁，针对绿光 \(550\) 纳米设计，膜厚约 \(100\) 纳米——大约一个病毒的大小。这层纳米薄膜让镜头反射率明显下降，透光更多、杂光更少；而它在白光下泛出的紫色，本章结尾的挑战题会让你自己推出来。
\end{mathbridge}

\begin{challenge}
这一章反复出现的 \(\lambda/a\)，其实是一张更大地图的一角。用傅里叶的语言说：一个孔径（缝、头发、瞳孔、望远镜口径）对光的限制越厉害，光被铺开的角度范围就越宽——“空间上的压缩”必然换来“方向上的摊开”，两者互为傅里叶变换。单缝越窄衍射越宽，和一段声音越短促其频率越不确定，是同一条数学定理的两张面孔。这条定理在量子篇会换一副庄严的面孔回来：单缝实验里光子穿过窄缝后方向的散开，正是位置与动量不确定关系（第 43 章）的预演——它不是仪器粗糙造成的，而是波的本性。

傅里叶光学还有一句漂亮的直觉：\emph{透镜能把方向变成位置}。平行光以某个角度射入透镜，会被聚到焦平面上的一个点；换一个入射方向，落点就换一个位置。也就是说，透镜把衍射图案里“向哪个方向散开”的信息，原封不动地翻译成“落在焦面上哪里”，等于在焦平面上实时画出孔径的傅里叶变换图。衍射远场、频谱分析、第 37 章要讲的成像理论，都从这一句长出来。

这张地图还顺手给所有光学仪器判了死刑：任何口径为 \(D\) 的圆孔，衍射都会把点光源摊成一个角半径约 \(1.22\lambda/D\) 的亮斑，两个点光源靠得比这个角度还近就无法分辨（瑞利判据）。代入人眼：瞳孔直径约 \(3\) 毫米，绿光 \(550\) 纳米，
\[
  \theta_{\min} \approx \frac{1.22\times 550\times10^{-9}\ \text{米}}{3\times10^{-3}\ \text{米}}
  \approx 2\times10^{-4}\ \text{弧度},
\]
约合一角分——正好对应视力表上 \(1.0\) 那一行。你的眼睛分辨率的极限，不是视网膜不够精密，而是光波在瞳孔门口的衍射。同理，芯片光刻机要用波长 \(13.5\) 纳米的极紫外光：想刻更细的电路，就得用更短的“光尺子”，这是衍射极限在半导体工业里值几百亿美元的那一面。
\end{challenge}

\step{模型的边界}

\begin{mainline}
波动模型本章战无不胜，但先给它划好国界，免得你在别处越界。

边界一：\emph{它依然不是光的最终面孔}。把光强调到极弱——暗到平均一次只有“一小份”光到达探测器——波动模型预言能量会像潮水一样连续均匀地铺到屏上，可探测器听到的却是一粒一粒离散的“咔哒”声。弱光下的双缝实验里，屏上先出现的是一个个离散的亮点，点越攒越多，才慢慢显出干涉条纹的形状。波动模型预言对了条纹的\emph{位置}，却说不清为什么是“一粒一粒”到达。这不是仪器粗糙，是光还有第三副面孔；那副面孔属于量子篇（第 42 章起），第 37 章会先把三副面孔的关系摆清楚。

边界二：\emph{波长不可忽略的场合才有波动光学}。障碍物比波长大成千上万倍时，衍射角 \(\lambda/a\) 小得测不出，波动模型和光线模型的预言完全相同——这时用光线模型即可，它更省事。这不是光“变回了”粒子，而是近似按时到岗。

边界三：\emph{干涉能否看见，受相干性管辖}。普通热光源的相干长度往往只有微米量级，路径差稍大一点条纹就洗掉了。阳光下的肥皂泡膜够薄（光程差在相干长度内），所以颜色看得见；同一片阳光照双缝实验却做不好，因为两缝到屏的路径差超过了波列的“记忆”。
\end{mainline}

再列三种本章特有的典型误用：

\begin{itemize}[leftmargin=1.8em]
  \item \emph{“暗条纹处的能量凭空消失了。”}错。干涉只是给能量重新分了房间：暗处少掉的，亮处加倍拿回来，屏幕上的总能量等于两束光之和。能量守恒从来没有被“抵消”违反——抵消的是\emph{振动}，不是能量。
  \item \emph{“衍射是光撞到孔的边缘被弹开、刮弯了。”}错。惠更斯原理里没有任何碰撞：是没被挡住的波前残部继续发出子波，自然铺满前方。把孔再开大十倍，边缘还是那圈边缘，衍射却消失了——可见关键不在边缘，而在孔径尺寸与波长之比。
  \item \emph{把波形图当成光的飞行轨迹。}教科书上那条上下起伏的蛇形线，画的是“电场在各处的强弱和方向”，不是光在空间走出一条波浪线。光（在均匀介质中）依然沿直线飞；波动的是电磁场这个量，不是飞行路线。同理，双缝干涉也不是“两束光撞在一起打架”，而是两列波各自到达屏上每一点后，在该点按节拍叠加。
\end{itemize}

\step{回头解释开场现象}

现在，用新模型逐一收账。

\textbf{肥皂泡的颜色。}皂膜是一层水薄膜。阳光照到膜上，一小部分在上表面反射，另一部分钻进膜里、在下表面反射后再钻出来。这两束反射光走过不同的路，光程差约为 \(2nt\)；对某一波长，若光程差恰好凑成相长条件，那种颜色就在反射光里被加强，其余颜色被削弱。膜的厚度因重力从上到下逐渐变化，于是“哪种颜色被加强”随高度排队——彩色横带出现了；风一吹，厚度抖动，花纹就流动。白光里什么波长都有，所以每一圈位置总能找到自己的颜色。最精彩的是膜顶那块\emph{发黑}的区域：膜薄到远小于波长时，光程差 \(2nt\) 趋向零，可上表面反射带半波损失、下表面不带，两束反射光天生反拍——膜薄到极致，反射光相消，膜在反射光里变成了黑色。一块透明的水膜，靠“光减光”变成了黑色。黑色不是脏，是干涉写给世界的签名。

马路上的油花是同一个原理，只是油膜厚度不均匀，彩色花纹更凌乱。大自然里这套把戏更是被用得炉火纯青：孔雀羽毛的蓝绿、蜻蜓翅膀闪现的金属光泽、一些甲虫壳的颜色，都不是颜料，而是纳米级的薄层结构玩出来的干涉色——这种颜色称为结构色。结构色有个颜料永远学不会的本事：换一个角度看，颜色就变，因为光程差随角度变了。

\textbf{光盘的彩虹。}已在光栅一节对账：\(d\sin\theta=m\lambda\)，\(1.6\) 微米的轨道间距把不同颜色掰向不同角度。镜子没有刻痕，没有彩虹；刻痕磨平，彩虹消失。

\textbf{指缝的条纹。}两指之间是一道窄缝，单缝衍射把灯光沿“缝被捏窄的方向”摊开成明暗条纹，角宽度约 \(2\lambda/a\)，缝越窄摊得越开——第一题的正确答案是（c），方向与直觉相反。眯起眼睛看路灯时睫毛之间的小缝，也会送你同样的光芒。

\textbf{双缝暗处变亮。}暗条纹中心是两列波反拍相消的位置：两路振动恰好互相抵消。挡住一条缝，只剩一列波独舞，无人与它抵消，该点反而分到单缝衍射的均匀亮度——所以第二题的答案也是（c）：光变少了，那个点却亮了。同时整屏的总能量确实减小了，只是重新分配后，原来的暗处分到了一杯羹。反常识的不是实验，是“光越多处处越亮”这句旧直觉。

\step{脑洞实验与挑战题}

\begin{thought}[用整个地球当一架望远镜]
瑞利判据说，想看更小的东西，要么缩短波长，要么把口径 \(D\) 做大。天文学家把这个逻辑推到了疯狂的极致：既然造不出地球那么大的望远镜镜片，那就让散布在全球的射电望远镜\emph{同时}观测同一个目标，各自精确记录到达的波，事后再把记录对齐节拍、做干涉合成——这正是本章“按节拍叠加”的工程版。2019 年，事件视界望远镜用分布在全球、基线长达上万公里的射电阵列，在 \(1.3\) 毫米波长上工作，等效口径接近地球直径，分辨率达到约二十微角秒——相当于从北京看清上海的一枚硬币——最终拍下了 M87 星系中心黑洞的阴影照片。那圈著名的橙色光环，本质上是人类用整个地球做孔径、用干涉原理“洗”出来的一张衍射极限照片。从肥皂泡到黑洞，中间只隔着一个“光程差”。
\end{thought}

\begin{thought}[一次只放一个光子]
把双缝实验的光源调暗、再调暗，直到平均而言同一时刻最多只有“一份”光在装置里飞行。波动模型预言条纹会连续变淡；实验看到的却是：屏上每次只亮起一个离散的点，落在哪里无法预言，但成千上万个点累积起来，排成的正是双缝干涉条纹——仿佛每一份光都“知道”两条缝的存在，各自按条纹的概率图选择落点。这还不是最怪的：若在缝旁放探测器记录“它到底走了哪条缝”，条纹立刻消失，屏上只剩两团几何亮斑。要说明白这一幕，波动和粒子这两副经典面孔都不够——它不是光在波和粒子之间来回变身，而是需要一套全新的语言：量子态的叠加与概率幅。这场“经典物理的灾难”将在第 42 章正式爆发，第 43 章起用量子力学重新审理。本章只需记住：双缝条纹这张“节拍账”，连一份一份到达的光也照记不误，这才是它最深的地方。
\end{thought}

\begin{puzzles}
\item 双缝实验中，暗条纹处两列波“互相抵消”。有人说：“两束光的能量在对消中被消灭了一部分，否则亮条纹凭什么比两束光简单相加还亮？”请找出这句话的错误。\emph{思路提示}：把明暗渐变的条纹亮度对整个屏幕求平均，与“两束光各自亮度之和”比较。抵消的是屏上\emph{该点}的振动，能量只是从暗处搬到了亮处，总账一分不差。
\item 用白光（阳光）做双缝实验，中央条纹是什么颜色？从中央往外数，第一圈彩色里，靠内的是红色还是紫色？\emph{思路提示}：\(\theta=0\) 处所有波长的路程差都是零，全部同拍。条纹间距 \(\Delta y=L\lambda/d\) 与波长成正比，想想红光和紫光谁的波长更长。
\item 你站在大楼背阴的一侧，看不见楼顶的灯，却能清楚收到调频广播；可见光和调频电波同为电磁波，待遇为何天差地别？\emph{思路提示}：衍射角约 \(\lambda/a\)。可见光波长不到一微米，调频广播波长约三米，而楼的尺寸是几十米。分别比较“波长与障碍物尺寸”这个比值，看谁的球面子波能铺满楼的阴影区。
\item 相机镜头的增透膜按绿光（\(550\) 纳米）设计，让绿光反射相消。在白光下斜看镜头，反射光常呈淡淡的紫色。为什么是紫色，而不是绿色？\emph{思路提示}：绿光被“消灭”得最干净，红光和蓝紫光只被部分消减。反射光里缺了绿色，剩下的颜色混合起来是什么？顺便想想：厂家为什么偏偏选绿光做设计波长——人眼对什么颜色最敏感？
\end{puzzles}

至此，光的三副面孔我们见到了两副：第 33 章的射线，本章的波。波的面孔还没看全——光波是横波，它的振动还有“朝向”这一维，那将引出偏振，以及光钻进物质内部之后的故事（第 36 章）。而两副面孔如何统一、第三副面孔为何必要，是第 37 章和量子篇的任务。一张光盘、一个肥皂泡、一道指缝，已经把“光到底是什么”这个问题，又往深处推了一大步。
